Remove the receipt from a delivery platform and every fraud control has to infer
instead of verify. We built Errandly's controls around that fact. Price claims are
bound to the lines the requester ordered and judged against a reference learned with
one vote per runner inside a human-set band, and runners who habitually claim just
under the limit lose their vote. Social closure discounts trust but never accuses; an
accusation needs money cycling among mutual friends. And because the dispatcher's own
preference for friends creates in-group circulation, ring flags are tested against an
acceptance model fitted to how students outside every suspected group actually
respond to offers.

In simulation the cycle search found rings that structure-only, money-only and
graph-mining baselines missed, and the policy test removed most of the false flags
that the router's preference for friends had caused. The test's first version failed
in a way we did not anticipate: by borrowing the router's own weights as its model of
behaviour, it cleared almost every coordinated ring. Fitting that model to logged
acceptances repaired most of the failure without the service cost of more
exploration. What no version can repair is the case in which honest friends and a
careful ring route errands to each other in the same way; there, only a person who
reads the errands can decide. The next step is a pilot with ethics approval, which
can replace our assumed behaviour and hand-set thresholds with measured ones.
